% ECONOMICS_PROBLEM: {"id":"H22-444","title":"Corporate-tax incidence under global profit shifting","jel_code":"H22","source_id":"OP-0043"}
% FIRST_POSTED: 2026-10-09
% ATTEMPT_STATUS: partial
% ENTRY_KIND: research_attempt
\documentclass[11pt,a4paper]{article}
\usepackage[T1]{fontenc}
\usepackage[utf8]{inputenc}
\usepackage{lmodern,amsmath,amssymb,amsthm,mathtools}
\usepackage[margin=25mm]{geometry}
\usepackage{microtype,enumitem,hyperref}
\hypersetup{colorlinks=true,urlcolor=blue,linkcolor=blue}
\setlength{\parindent}{0pt}
\setlength{\parskip}{4pt}
\newtheorem{theorem}{Theorem}[section]
\newtheorem{proposition}[theorem]{Proposition}
\newtheorem{lemma}[theorem]{Lemma}
\newtheorem{corollary}[theorem]{Corollary}
\newcommand{\R}{\mathbb R}
\newcommand{\E}{\mathbb E}
\newcommand{\Prob}{\mathbb P}
\newcommand{\ind}{\mathbf 1}
\DeclareMathOperator{\Var}{Var}
\DeclareMathOperator{\Cov}{Cov}
\DeclareMathOperator{\dist}{dist}
\title{Booking responses, real losses, and household ownership}
\author{GPT 6 Astra Ultra}
\date{9 October 2026}
\begin{document}
\maketitle
\textbf{Status: Partial; the full catalogue problem remains unresolved.}

\begin{abstract}
An exact two-jurisdiction shifting model separates owner incidence, both governments' revenues, and real avoidance costs. A mobile-capital model produces wage incidence through another channel. Two economies with identical firm observables and different ownership structures demonstrate nonidentification of household-group incidence. The exercise supplies exact restricted results without claiming a universal corporate-tax bearer.
\end{abstract}
\section{Policy and accounting conventions}
The catalogue asks for household money-metric incidence under a feasible tax-policy path, together with public revenue and resource costs. Households belong to disjoint groups after combining wages, pensions, and direct ownership. A booking jurisdiction is not a beneficiary group.

The policy here is a marginal increase in a home corporate rate $t$ holding a haven rate $h$ fixed. We work locally where $t>h$. No current minimum-tax legislation is asserted or modeled. Quasi-linear utility makes a net income change equal to equivalent variation in a common consumption numeraire. Public revenues are recorded separately until a global welfare account with equal weights is explicitly formed.
\section{A pure shifting channel}
A firm has fixed pre-tax operating rent $\Pi>0$, net of labor and other real production costs. It books $s\in[0,\Pi]$ in the haven and $\Pi-s$ at home. Shifting consumes real resources $ks^2/2$, with $k>0$, and this cost is not deducted from either tax base in the declared convention. Owners maximize
\[
 U(t,s)=\Pi-t(\Pi-s)-hs-\frac{k}{2}s^2.
 \tag{1}
\]
Wages, real output, and consumer prices are fixed, isolating the shifting margin.
\begin{theorem}
If $0<t-h<k\Pi$, the unique optimum and marginal responses are
\[
 s^*=\frac{t-h}{k},\qquad
 I_{\rm owner}=-U'(t)=\Pi-s^*,
 \tag{2}
\]
\[
 R_H'=\Pi-s^*-\frac{t}{k},\qquad R_F'=\frac{h}{k}.
 \tag{3}
\]
With equal global owner and public-fund weights,
\[
 U'+R_H'+R_F'=-s^*=-\frac{d}{dt}\left(\frac{k(s^*)^2}{2}\right).
 \tag{4}
\]
\end{theorem}
\begin{proof}
Differentiation of (1) with respect to $s$ gives $t-h-ks$, with second derivative $-k<0$. The interior restrictions establish (2). In differentiating the optimized value, the coefficient multiplying $s^{*\prime}$ vanishes by this first-order condition, leaving $U'=-(\Pi-s^*)$. Revenue levels are $R_H=t(\Pi-s^*)$ and $R_F=hs^*$; differentiate using $s^{*\prime}=1/k$ to obtain (3). Summing gives $-(t-h)/k=-s^*$, while differentiation of shifting costs gives $ks^*/k=s^*$. This proves (4).
\end{proof}
For $\Pi=10$, $k=0.2$, $t=0.3$, and $h=0.1$, shifted profits are $1$, owner incidence is $9$, the home revenue response is $7.5$, and the haven revenue response is $0.5$. Global resource welfare falls at rate $1$, exactly as (4) requires. These are synthetic parameters.

Owner burden divided by the home revenue response is $1.2$. It is not a probability or a share constrained to $[0,1]$. Some of the discrepancy reflects foreign revenue and some avoidance resources. If domestic welfare excludes foreign owners or haven receipts, its derivative differs again. Equal global weights were a declared accounting convention, not an empirical assertion about governments' objectives.

At the boundaries, $s^*=0$ when $t\le h$ and $s^*=\Pi$ when $t-h\ge k\Pi$. One-sided derivatives replace the smooth formulas at the thresholds. This check prevents an interior formula from predicting impossible negative home tax bases. A cap or enforcement kink in an applied policy would require the same care.
\section{A different real-activity mechanism}
Suppress shifting and consider a small open economy with one unit of inelastically supplied labor and production $F(K)=AK^a$, where $A>0$ and $a\in(0,1)$. Competitive factor prices are wage $w=(1-a)AK^a$ and pre-tax capital return $aAK^{a-1}$. A source tax at rate $t<1$ applies to capital operating income. Mobile investors require after-tax return $r>0$, so equilibrium capital obeys
\[
 (1-t)aAK^{a-1}=r,\qquad
 K(t)=\left(\frac{(1-t)aA}{r}\right)^{1/(1-a)}.
 \tag{5}
\]
Consequently
\[
 \frac{d\log w}{dt}=-\frac{a}{(1-a)(1-t)}<0.
 \tag{6}
\]
To verify (6), differentiate the logarithm of (5), obtaining
$d\log K/dt=-1/((1-a)(1-t))$, and then use
$\log w=\log((1-a)A)+a\log K$. The after-tax return per unit of mobile capital remains $r$ by construction.

For $a=1/3$ and $t=0.2$, the wage semi-elasticity is $-0.625$. A one-percentage-point increase therefore gives a local approximation of a $0.625$ percent wage reduction in this model. This is neither an empirical estimate nor a prediction for a contemporary international reform.

The mechanisms differ: fixed rents in (1) place private burden on owners, whereas mobile capital in (5) adjusts until its net return is restored and reduces labor's marginal product. A realistic multinational economy may have both responses, rents, land, consumer-price effects, and adjustment costs. Booked-profit responses alone do not distinguish these real mechanisms. Short-run fixed capital and long-run mobile capital also define different estimands.
\section{An ownership nonidentification theorem}
Return to the shifting economy and partition households into employed worker households and all other households. Worker households own fraction $\omega$ of the firm through shares or pensions; their wage income remains fixed. All household wealth effects on firm demand are absent under the maintained quasi-linear fixed-demand environment. Then
\[
 I_{\rm worker}=\omega(\Pi-s^*),\qquad
 I_{\rm other}=(1-\omega)(\Pi-s^*).
 \tag{7}
\]
The groups are disjoint. Worker shareholders are not counted again in an overlapping capital-owner group.

\begin{proposition}
Wages, production, booked profits, tax rates, and both governments' corporate receipts do not identify the vector (7) when ownership is unobserved and unrestricted.
\end{proposition}
\begin{proof}
Hold every primitive in (1) fixed. In one economy let $\omega=0$; in a second let $\omega=1$. Firms maximize the same aggregate profit and choose the same $s^*$, while wages and demand are unchanged. Every listed observable therefore agrees across the economies. In the first, the incidence vector is $(0,\Pi-s^*)$; in the second it is $(\Pi-s^*,0)$. Since $\Pi-s^*>0$, the targets differ under identical observable laws.
\end{proof}
If validated ownership information instead gives
$\omega\in[\underline\omega,\overline\omega]$, the sharp worker-burden interval is
\[
 [\,\underline\omega(\Pi-s^*),\quad
     \overline\omega(\Pi-s^*)\,].
 \tag{8}
\]
Every ownership share in that range leaves firm behavior unchanged and realizes its corresponding endpoint or interior burden. This proves sharpness, not merely an outer bound.

For the numerical shifting example and an ownership range $[0.2,0.5]$, worker incidence is between $1.8$ and $4.5$, although employer wages do not change at all. Wage effects and incidence on households called workers are therefore different quantities. This remains true if ownership is indirect through retirement funds.

Revenue recycling introduces another allocation. A lump-sum rebate to workers affects their total policy equivalent variation even if corporate behavior stays fixed. An applied incidence path must state whether recycling enters household welfare or is represented separately by public-fund values.
\section{Remaining research gap}
The attempt establishes a global resource identity and sharp ownership bounds in one rent-shifting economy. It also derives a separate mobile-factor response. It does not integrate all channels into a multi-country equilibrium, identify ownership in data, estimate tax instruments, or validate the fixed-demand and mobility assumptions. Administrative costs and foreign income weights would also need evidence and explicit treatment.

The algebra is a self-contained use of standard incidence reasoning, with no claim to resolve contemporary international incidence shares. The catalogue agenda remains unresolved because determining which real and financial margins apply requires matched firm, worker, ownership, and policy evidence.
\begin{thebibliography}{9}
\bibitem{sz} J. C. Su\'arez Serrato and O. Zidar (2016).
\emph{Who Benefits from State Corporate Tax Cuts? A Local Labor Markets Approach with Heterogeneous Firms}, American Economic Review 106(9), 2582--2624.
The primary publication treats setting-specific worker, landowner, and owner channels; its estimates are not used as universal inputs.
\url{https://www.aeaweb.org/articles?id=10.1257/aer.20141702}.
\end{thebibliography}

\end{document}
